\chapter{Deployment and Web}

\section{Key Internet / Web Concepts}
\dfn{Protocols}{
HTTP (hypertext), TCP (reliable transport), IP (addressing/routing), Ethernet (link).
}
\dfn{DHCP}{
Dynamic Host Configuration Protocol; assigns IPs automatically.
}
\dfn{IP Addressing}{
Private vs public IP; servers reachable via public IP or domain.
}
\dfn{DNS}{
Resolves domain names to IP addresses.
}
\dfn{HTTP Request--Response}{
Client sends request, server executes logic/DB access, returns response (HTML/CSS/JS/media), browser renders.
}
\dfn{Browser}{
Issues HTTP, interprets HTML/CSS/JS, renders DOM; offers dev tools.
}
\dfn{Stacks}{
XAMPP (Apache, MariaDB/MySQL, PHP/Perl); MEAN (Mongo, Express, Angular, Node); .NET; Java EE; growing use of Python frameworks.
}

\section{Web Client Fundamentals}
\dfn{HTML}{
Markup language composed of tags and text; DOCTYPE + head + body.
}
\dfn{CSS}{
Styles via rules (selector \{ attr: value; \}); inline, head, or external stylesheet.
}
\dfn{JavaScript}{
Lightweight, interpreted, prototype-based language for client-side interactivity (also Node.js server-side).
}

\subsection*{HTML Skeleton}
\begin{lstlisting}
<!DOCTYPE html>
<html lang="en">
  <head>
    <meta charset="UTF-8">
    <title>Demo</title>
    <link rel="stylesheet" href="style.css">
  </head>
  <body>
    <h1>Hello Web</h1>
  </body>
</html>
\end{lstlisting}

\subsection*{Form Example}
\begin{lstlisting}
<form method="post" action="/submit">
  <input type="text" name="city" value="Lisbon" size="15" maxlength="15">
  <input type="submit" value="Send">
</form>
\end{lstlisting}

\subsection*{CSS Examples}
\begin{lstlisting}
/* external style.css */
body { background: yellow; }
p    { color: blue; }
\end{lstlisting}

\subsection*{JS Example}
\begin{lstlisting}
<!DOCTYPE html>
<html>
<head>
  <script>
    function teste(){ alert("olá"); }
  </script>
</head>
<body>
  <input type="button" value="Olá" onclick="teste();">
</body>
</html>
\end{lstlisting}

\section{Python Web with Flask}
\dfn{Flask}{
Lightweight WSGI microframework using Werkzeug (WSGI toolkit) and Jinja2 (templating).
}
\dfn{Routing}{
Bind URL paths to view functions with \texttt{@app.route}.
}
\dfn{Static Files}{
Served from \texttt{/static}; reference via \texttt{url\_for('static', filename='style.css')}.
}

\dfn{Templates}{
Jinja2 templates rendered with \texttt{render\_template}.
}

\subsection*{Minimal App}
\begin{lstlisting}
from flask import Flask
app = Flask(__name__)

@app.route("/")
def index():
    return "Hello, World!"

if __name__ == "__main__":
    app.run()
\end{lstlisting}

\subsection*{Routing Example}
\begin{lstlisting}
@app.route("/")
def index():
    return "Index Page"

@app.route("/hello")
def hello():
    return "Hello, World"
\end{lstlisting}

\subsection*{Template Rendering}
\begin{lstlisting}
from flask import Flask, render_template
app = Flask(__name__)

@app.route("/hello/")
@app.route("/hello/<name>")
def hello(name=None):
    return render_template("hello.html", name=name)
\end{lstlisting}

Template \texttt{templates/hello.html}:
\begin{lstlisting}
<!doctype html>
<title>Hello from Flask</title>
{% if name %}
  <h1>Hello {{ name }}!</h1>
{% else %}
  <h1>Hello, World!</h1>
{% endif %}
\end{lstlisting}

\subsection*{Simple Form Handling}
\begin{lstlisting}
from flask import Flask, request, render_template
app = Flask(__name__)

@app.route("/", methods=["GET", "POST"])
def form():
    if request.method == "POST":
        author = request.form.get("Author")
        phrase = request.form.get("phrase")
        # save/process...
        return f"Saved: {author} - {phrase}"
    return render_template("form.html")
\end{lstlisting}

Form template \texttt{templates/form.html}:
\begin{lstlisting}
<form action="/" method="POST">
  <p>author <input type="text" name="Author"></p>
  <p>phrase <input type="text" name="phrase"></p>
  <p><input type="submit" value="Save"></p>
</form>
\end{lstlisting}

\subsection*{Serving Static Assets}
\begin{lstlisting}
<!-- in template -->
<link rel="stylesheet" href="{{ url_for('static', filename='style.css') }}">
\end{lstlisting}

\section{Deployment Notes}
\begin{itemize}
    \item Run locally: \texttt{python app.py} then visit \texttt{http://localhost:5000}.
    \item Production: use WSGI server (gunicorn/uwsgi) + reverse proxy (nginx); set \texttt{FLASK\_ENV=production}.
    \item Static vs dynamic content: CDN/static folder for assets; keep secrets out of repo; configure environment variables.
    \item For quick demos: platforms like Render, Railway, or Heroku-like PaaS; ensure requirements.txt and Procfile.
\end{itemize}

\section{Practice / Exam Prompts}
\pr{HTTP flow}{
Describe the path from typing a URL to page render: DNS lookup, TCP three-way handshake, HTTP request, server logic/DB, HTTP response, browser parse/render.
}
\pr{Flask routing}{
Add a new route \texttt{/about} returning HTML; show how to link to it from a template via \texttt{url\_for}.
}
\pr{Forms}{
Explain POST vs GET for forms; how to read \texttt{request.form} and validate input.
}
\pr{Static files}{
Where to place CSS/JS in Flask and how to reference them with \texttt{url\_for}.
}
\pr{Deployment checklist}{
Use gunicorn behind nginx, set \texttt{DEBUG=False}, configure env vars for secrets, enable logging, and serve static via nginx/CDN.
}
